\documentclass{article}

\usepackage{amsmath,amssymb}
\usepackage{xurl}
\usepackage[hidelinks]{hyperref}

% Shared commands for notes in docs/paper-gaps/.

\DeclareUnicodeCharacter{2080}{\ifmmode{}_{0}\else\textsubscript{0}\fi}
\DeclareUnicodeCharacter{2081}{\ifmmode{}_{1}\else\textsubscript{1}\fi}
\DeclareUnicodeCharacter{2082}{\ifmmode{}_{2}\else\textsubscript{2}\fi}
\DeclareUnicodeCharacter{2099}{\ifmmode{}_{n}\else\textsubscript{n}\fi}

\newcommand{\Lean}{\textsc{Lean}}
\ExplSyntaxOn
\NewDocumentCommand{\leanid}{m}{
  \ifmmode
    \textsf{\detokenize{#1}}
  \else
    \group_begin:
    \urlstyle{sf}
    \tl_set:Nn \l_tmpa_tl {#1}
    \tl_replace_all:Nnn \l_tmpa_tl {\_} {_}
    \exp_args:NV \path \l_tmpa_tl
    \group_end:
  \fi
}
\ExplSyntaxOff
\providecommand{\tr}{\operatorname{tr}}
\newcommand{\tnleanIssueBase}{https://github.com/LionSR/TNLean/issues/}
\newcommand{\tnleanPrBase}{https://github.com/LionSR/TNLean/pull/}
\newcommand{\tnleanDocBase}{https://sirui-lu.com/TNLean/docs/}
\newcommand{\ghissue}[1]{\href{\tnleanIssueBase#1}{GitHub issue \##1}}
\newcommand{\ghpr}[1]{\href{\tnleanPrBase#1}{PR \##1}}
\newcommand{\doclink}[1]{\href{\tnleanDocBase#1.html}{\path{#1}}}

% Dirac notation and the identity operator, as in blueprint/src/macros/common.tex.
% \providecommand keeps note-local definitions authoritative.
\usepackage{dsfont}
\providecommand{\ket}[1]{|#1\rangle}
\providecommand{\bra}[1]{\langle #1|}
\providecommand{\braket}[2]{\langle #1 | #2 \rangle}
\providecommand{\ketbra}[2]{|#1\rangle\!\langle #2|}
\providecommand{\Id}{\mathds{1}}
\providecommand{\one}{\mathds{1}}
\providecommand{\C}{\mathbb{C}}
\providecommand{\MN}[1]{M_{#1}(\C)}
\providecommand{\MD}{M_D(\C)}

% External referencing for paper sources.  The build helper writes sanitized
% auxiliary files under build/paper-aux/ and excludes duplicated source labels.
\usepackage{xr}

% Import a generated paper auxiliary file with a caller-supplied label prefix.
% The candidate paths support builds from docs/paper-gaps/, the repository root,
% and build/paper-aux/ itself.
\newcommand{\importpaperaux}[2]{%
  \IfFileExists{../../build/paper-aux/#2.aux}{%
    \externaldocument[#1]{../../build/paper-aux/#2}%
  }{%
    \IfFileExists{build/paper-aux/#2.aux}{%
      \externaldocument[#1]{build/paper-aux/#2}%
    }{%
      \IfFileExists{#2.aux}{%
        \externaldocument[#1]{#2}%
      }{}%
    }%
  }%
}

% General external reference macros
\newcommand{\paperref}[2]{\ref{#1-#2}}
\newcommand{\papereqref}[2]{(\ref{#1-#2})}

% CPSV16 source (1606.00608 / MPDO-22-12-17-2.tex) external document and shortcuts
\importpaperaux{cpsv16-}{1606.00608/MPDO-22-12-17-2-tnlean}
\newcommand{\cpsvref}[1]{\paperref{cpsv16}{#1}}
\newcommand{\cpsveqref}[1]{\papereqref{cpsv16}{#1}}

% Verdict marker: \gapnote{<kind>}{<status>}. Kind is the policy
% classification (clarification, local-correction, scope-restriction,
% unfaithful, false-source, open-gap); status is open, wip (elimination
% underway), resolved, or historical. Machine-read by the site index;
% prints a verdict line.
\newcommand{\gapnote}[2]{\par\noindent\textbf{Verdict:} #1 (#2).\par}


\title{The Endpoint Rank-Product Exponent in MPU Index Well-Definedness}
\author{}
\date{2026-08-12}

\begin{document}
\maketitle

\section{The printed exponent}

In the proof of Proposition~IV.2 of \cite{Cirac2017MPU}, the right and left
source-cut ranks of the $k$-site blocked tensor are printed as
\[
  r_k\ell_k=d^k.
\]
The blocked physical dimension is $d^k$.

\section{Consistent rank-product identity}

Theorem~III.8 states that for a simple tensor of physical dimension $q$, the
corresponding ranks obey
\[
  r\ell=q^2.
\]
Applying this statement to the $k$-site block, with $q=d^k$, gives
\[
  r_k\ell_k=(d^k)^2=d^{2k}.
\]
The same square occurs in the index discussion and in the later blocking
argument. Thus the occurrence of $d^k$ in the proof of Proposition~IV.2 is a
typographical exponent error; the consistent endpoint identity is $d^{2k}$.

\section{Formal treatment}

The simultaneous saturation theorem
\leanid{MPOTensor.blockingRanks_eq_pow_mul_of_products} does not prove this
rank-product identity. It takes the two endpoint identities
\[
  r_{k_0}\ell_{k_0}=d^{2k_0},
  \qquad
  r_k\ell_k=d^{2k}
\]
as explicit hypotheses and combines them with the independently proved
blocking upper bounds. Its focused accessors are
\leanid{MPOTensor.rightRank_blockTensor_eq_pow_mul_of_products} and
\leanid{MPOTensor.leftRank_blockTensor_eq_pow_mul_of_products}. A later
application may obtain the endpoint identities from the formal counterpart of
Theorem~III.8 under its own assumptions.

\section{Verdict}

\textbf{Verdict: local correction (resolved).}  The exponent in the proof of
Proposition~IV.2 is corrected from $d^k$ to $d^{2k}$. This does not add an
assumption beyond the rank-product identity used by that proof, but the formal
saturation result keeps that identity explicit rather than claiming it from
blocking alone.

\bibliographystyle{alpha}
\bibliography{references}

\end{document}
